\begin{ledgroupsized}[r]{120mm}
\footnotesize
\pstart
\noindent\textbf{Überlieferung:}
\pend
\end{ledgroupsized}
\begin{ledgroupsized}[r]{114mm}
\footnotesize
\pstart \parindent -6mm
\makebox[6mm][l]{\textit{L}}%
Konzept:
LH~XXXVII~5 Bl.~189. 
Ein Zettel (ca.~13~x~15~cm.) mit einem überstehenden Streifen unten rechts (ca.~7~x~3,5~cm.);
\href{\Scan{LH_37_05_189wz}}%
{Wasserzeichenfragment} am oberen Rand;
Papier\-erhaltungsmaßnahmen;
alle Ränder beschnitten.
Eine Seite auf Bl.~189~r\textsuperscript{o};
Bl.~\href{\Scan{LH_37_05_189v}}{189~v\textsuperscript{o}} 
überliefert einen eigenhändigen Briefentwurf:
\selectlanguage{french}%
\textit{Monsieur} \lbrack/\rbrack\
\textit{je vous demande pardon de ce que la premiere que j'ay} \lbrack/\rbrack\ \textit{l'honneur de vous ecrire est si precipitée. C'est pour profiter} \lbrack/\rbrack\ \textit{de l'occasion de vous envoyer ce qui y est joint. Je vous écriray} \lbrack/\rbrack\ \textit{plus amplement par la poste prochaine. Il me déplaist de ne} \lbrack/\rbrack\ \textit{vous pouvoir point envoyer la suite entiere des journaux de Paris.} \lbrack/\rbrack\ \textit{Comme je n'en reçois \sout{qu' apres avoir} que ce qui a passé par plusieurs mains à la} \lbrack/\rbrack\ \textit{cour, il s'en perd beaucoup, et quelques fois on ne me porte que des} \lbrack/\rbrack\ \textit{lambeaux. Comme en effect un des trois que voicy est imparfait.}
\pend
\end{ledgroupsized}
\vspace{5mm}
\begin{ledgroup}
\footnotesize
\pstart
\noindent%
\textbf{Datierungsgründe:}
Das Wasserzeichen im Träger von N.\;\ref{60356} ist für den Zeitraum Oktober 1690 bis Oktober 1697 mehrfach belegt.
Der Adressat des Entwurfs auf Bl.~189~v\textsuperscript{o}, offenbar als Begleitbrief für die Zusendung dreier Zeitschriften gedacht,
konnte noch nicht ermittelt werden.
\pend
\pstart
Am Ende von N.\;\ref{60356} (\refpassage{37_05_189_2a}{37_05_189_2b}) beruft sich Leibniz auf gewisse Ergebnisse der \glqq Dynamica\grqq, womit entweder bestimmte Schriften mit dem Titel "dynamica" (etwa die Entwürfe zur \cite{01354}\textit{Dynamica de potentia} von 1689, das \cite{02167}\textit{Essay de Dynamique} von 1692 oder aber das \cite{02032}\textit{Specimen Dynamicum} von 1695) gemeint sein könnten,
oder vielmehr die Disziplin der Dynamik im Allgemeinen, die der "Phorographica" folgen soll.
In diesem Zusammenhang siehe das vollständigste unter den erhaltenen Inhaltsverzeichnissen zur \textit{Dynamica de potentia}: Eine Randanmerkung von Leibniz zum allgemeinen Teil, insbesondere zur "Sect.\;2. De Motu in Genere", lautet "Addenda esset sectio phorographica" (\cite{01354}\textit{CFP} I, S.\;5 [A]).
\pend
\end{ledgroup}
\selectlanguage{latin}
\vspace{8mm}
\pstart%
\normalsize%
\noindent%
\href{\Scan{LH_37_05_189r}}%
{\lbrack189~r\textsuperscript{o}\rbrack}%
\edlabel{37_05_189_1a} %
\edtext{}{%
{\xxref%
{37_05_189_1a}{37_05_189_1b}}%
\lemma{\lbrack\textit{Fig.~1}\rbrack}%
\Bfootnote{%
\textit{(1)}~$\ind{1}C\ind{2}C=$
\textit{(a)}~\textit{\ind{1}A\ind{2}A}
\textit{(b)}~$\ind{2}A\ind{1}A-\ind{1}A\ind{1}C=
\ind{2}B\ind{2}B$\lbrack\textit{!}\rbrack\ $-$
\textit{(2)}~\bars{(1mo) \textit{erg.}}
$\ind{1}A\ind{1}C+\ind{1}C\ind{1}B=\ind{1}A\ind{2}A-\ind{2}B\ind{1}B=$
\bars{(2do) \textit{erg.}} 
$\ind{2}A\ind{3}C+\ind{3}C\ind{3}B=$
\textit{(a)}~$\ind{3}A\ind{2}A-\ind{2}B$
\textit{(b)}~\bars{(3tio) $\ind{3}B\ind{2}B-\ind{2}A\ind{3}A$ $=$ (4to) \textit{AB} \textit{streicht Hrsg.}}
\textit{(c)}~(3tio) \textit{\ind{1}C\ind{2}C} (4to)
$\ind{1}C\ind{2}C=\ind{2}C\ind{3}C=$ (5to)
$\ind{2}A\ind{1}A-\ind{1}A\ind{1}C$ (6to) $=$
\textit{(3)}~\bars{(1mo) \textit{erg.}}
$AB=A(A)-B(B)$. $C(C)=A(A)-$
\textit{(4)}~\textls{Phorographica} \lbrack/\rbrack\
\bars{Posito ante et post ictum
\textit{(a)}~omnes
\textit{(b)}~ambo tendere quo centrum
\bars{\textit{C} \textit{erg.}}
et \textit{A} sequi \textit{B} praecedere
ante ictum \textit{erg.}}
(1mo)~\textit{L}}}
\pend 
\begin{edfigure}[0.5em]%
\includegraphics%
[scale=1]{gesamttex/edit_VIII-4/images/LH_37_05_189_d_189r.pdf} %
\edcaption[0.25em]{\lbrack\textit{Fig.~1}\rbrack}
\end{edfigure}
\pstart\centering	%
\edtext{%
\textls{Phorographica}\indexS{Phorographica}%
}{\lemma{}\Afootnote{%
\textit{Am oberen rechten Rand}: 
$\ind{1}A\ind{1}C+\ind{1}B\ind{1}C=\ind{1}A\ind{1}C+\ind{1}B\ind{1}C=+\ind{1}A\ind{2}A-\ind{1}B\ind{2}B$}}
\pend
\nobreak\vspace{0.5\baselineskip}
\pstart\noindent
Posito ante et post ictum\indexS{ictus} ambo tendere quo centrum\indexS{centrum gravitatis} \textit{C} et \textit{A} sequi \textit{B} praecedere\lbrack,\rbrack\
ante ictum\indexS{ictus}
(1mo)%
\edlabel{37_05_189_1b}
$\ind{1}A\ind{1}B=\ind{1}A\ind{2}A-\ind{1}B\ind{2}B=$ 
(2do) $\ind{1}A\ind{1}C+\ind{1}B\ind{1}C$%
\pend
\pstart\noindent
(3tio) $\ind{1}C\ind{2}C=\ind{1}A\ind{2}A-\ind{1}A\ind{1}C=$ (4to) $\ind{1}B\ind{2}B+\ind{1}B\ind{1}C$
\pend
\pstart\noindent
post ictum (5to) $\ind{3}A\ind{3}B=-\ind{2}A\ind{3}A+\ind{2}B\ind{3}B=$ (6to) $\ind{3}A\ind{3}C+\ind{3}B\ind{3}C$
\pend
\pstart\noindent
(7mo) $\ind{2}C\ind{3}C=\ind{2}A\ind{3}A+\ind{3}A\ind{3}C=$ (8vo) $\ind{2}B\ind{3}B-\ind{3}B\ind{3}C$%
\pend
\pstart\noindent
Hinc generaliter (9no) $AB=\PM\, A(A) \MP B(B)=$ (10mo) $AC+BC$ et 
\edtext{(11mo)}{%
\lemma{}%
\Bfootnote{%
(11mo) %
\textit{erg.~L}%
}}
$C(C)=A(A) \MP AC=$ 
(12mo) $B(B) \PM BC$. 
Ubi notandum aeqq.~2, 6, 
\edtext{\lbrack11\rbrack,}{%
\lemma{}%
\Bfootnote{%
12 %
\textit{L ändert Hrsg.}%
}}
coincidere respective cum
\edtext{4, 8, 12. Et}{\lemma{4, 8, 12.}\Bfootnote{%
\textit{(1)}~Et aeqq.~9. 
\textit{(2)}~Aequ %
\textit{(3)}~Et~\textit{L}}}
aeqq.~9, 10, 11, 12, fieri respective ex collatis 1 et 5, 2 et 6, 3 et 7, 4 et 8 et in aeqq. dictis 4 
\edtext{postremis anteriorem}{\lemma{postremis}\Bfootnote{%
\textit{(1)}~prius signum ante %
\textit{(2)}~in signo ambiguo explicari de statu ante ictum %
\textit{(3)}~anteriorem~\textit{L}}}
explicationem signi ambigui\indexS{signum ambiguum} explicari de statu ante ictum,\indexS{status ante ictum}
\edtext{posteriorem de statu}{\lemma{posteriorem de}\Bfootnote{%
\textit{(1)}~signo %
\textit{(2)}~statu~\textit{L}}}
post ictum,\indexS{status post ictum} et quae vera sunt de \textit{A} ante ictum,\indexS{ictus} vera esse de \textit{B} post ictum\indexS{ictus} et vice versa.
\edtext{Atque haec vera}{\lemma{Atque}\Bfootnote{%
\textit{(1)}~hoc verum est %
\textit{(2)}~haec vera~\textit{L}}} 
sunt etiamsi status non sumantur immediate 
\edtext{ante ictum\indexS{ictus} vel}{\lemma{ante ictum}\Bfootnote{%
\textit{(1)}~sed %
\textit{(2)}~vel~\textit{L}}} 
immediate post ictum\lbrack,\rbrack\indexS{ictus}
\edtext{id est quorsumcunque}{\lemma{id est}\Bfootnote{%
\textit{(1)}~sive coincidant %
\textit{(2)}~quorsumcunque~\textit{L}}} 
cadant \textit{\ind{2}A}, \textit{\ind{2}B}, \textit{\ind{2}C}. 
Et licet fingas differre ea quae sunt primi (casus) cum iis quae sunt secundo, et ita 
\edtext{sumas 3 pro secundo 2, et 4}{\lemma{sumas}\Bfootnote{%
\textit{(1)}~\textit{\ind{1}A}, \textit{\ind{2}A}, \textit{\ind{3}A}, \textit{\ind{4}A}, seu
\textit{(2)}~3 pro secundo 2,
\textit{(a)}~cum hic sumamus \textit{\ind{1}A}, \textit{\ind{2}A}, \textit{\ind{2}A}, \textit{\ind{3}A}
\textit{(b)}~et 4~\textit{L}}}
pro 3; sive \textit{\ind{1}A}, \textit{\ind{2}A}, \textit{\ind{3}A}, \textit{\ind{4}A}, pro \textit{\ind{1}A}, \textit{\ind{2}A}, \textit{\ind{2}A}, \textit{\ind{3}A}. 
\pend
\pstart
\edlabel{37_05_189_2a}%
\edtext{}{%
\xxref{37_05_189_2a}{37_05_189_2b}%
\lemma{Sequuntur \bracks{...} generalis}%
\Cfootnote{%
Gemeint ist wohl die dynamische Disziplin im Allgemeinen, der eine \textit{phorographica} vorausgehen soll, und nicht etwa eine bestimmte Schrift. Ein ähnlich formuliertes Prinzip findet sich dennoch in den Entwürfen zur \textit{Dynamica de potentia} von 1689: pars II, sect. 6 (\textit{De concursu corporum}), prop.\;14 bzw. 19 (\cite{01354}\textit{CFP} I, S.\; 413 bzw. 430 [A]).}}%
Sequuntur \indexS{Dynamica}\textls{Dynamica}, ubi haec regula
\edtext{generalis.\edlabel{37_05_189_2b} \textls{Quicquid} reperitur}{\lemma{generalis.}\Bfootnote{%
\textit{(1)}~Omnis veritas quae reperitur %
\textit{(2)}~\textls{Quicquid} reperitur~\textit{L}}} 
commune duobus quibusvis statibus ante ictum,\indexS{status ante ictum} et duobus quibusvis statibus post ictum\lbrack,\rbrack\indexS{status post ictum} id etiam commune 
\edtext{reperitur cuivis statui ante ictum\indexS{status ante ictum}}{\lemma{reperitur}\Bfootnote{%
\textit{(1)}~statui ante ic %
\textit{(2)}~cuivis statui ante ictum~\textit{L}}} 
cum quovis statu post ictum.\indexS{status post ictum} \pend
